% Vietnamese translation for OI Public Library
% Source-File: IOI中国国家候选队论文/2008/Day1/6.肖汉骏《例谈信息学竞赛分析中的“深”与“广”》/HunkShaw.ppt
\documentclass[11pt]{article}
\usepackage{oipl-vn}

\newcommand{\Oh}{\mathcal{O}}

\title{Bản trình chiếu: Phân tích sâu và rộng trong thi tin học}
\author{Xiao Hanjun}
\date{2008}

\begin{document}
\maketitle

\origtitle{Depth and breadth in informatics analysis}
\sourcepath{IOI China candidate team papers / 2008 / Day 1 / Xiao Hanjun.ppt}

\section{Phạm vi bản dịch}

Bản trình chiếu của Xiao Hanjun đi cùng bài nói ``Đường rẽ còn dài, trên dưới
kiếm tìm'', bàn về ``chiều sâu'' và ``độ rộng'' trong phân tích bài toán tin
học. Tệp PPT gốc trích được ít chữ hơn hai deck còn lại, nhưng các công thức
và đề mục cho thấy nó tập trung vào một ví dụ: sắp thứ tự đọc các bài viết
theo thời gian \(T_i\), giá trị \(V_i\), và ràng buộc thứ tự trong từng tác
giả.

Bản ghi chú này giữ lại mạch của slide: bắt đầu từ trường hợp đơn giản, dùng
phân tích đổi chỗ để tìm quy luật, rồi mở rộng góc nhìn bằng hình học và bao
lồi.

\section{Sâu và rộng}

``Sâu'' trong bài nói không có nghĩa là dùng kỹ thuật khó ngay từ đầu. Nó là
cách đi từng tầng: hạ bài toán về trường hợp đơn giản, tìm quan hệ thật sự
chi phối nghiệm tối ưu, rồi đưa quan hệ ấy trở lại bài toán đầy đủ. Quá trình
đó cần liên tục; mỗi bước phải kế thừa điều vừa chứng minh ở bước trước.

``Rộng'' là khả năng đổi góc nhìn. Một bài toán có thể được nhìn như sắp xếp,
quy hoạch động, hình học, đồ thị hoặc bất đẳng thức. Khi một biểu diễn chưa
đủ sắc, ta thử biểu diễn khác nhưng vẫn giữ liên hệ với kết luận đã có. Trong
ví dụ của slide, biểu thức tỉ số \(V_i/T_i\) dẫn tự nhiên tới cách nhìn hệ số
góc của vectơ.

\section{Bài toán thứ tự đọc}

Có \(N\) bài viết. Bài \(i\) cần thời gian đọc \(T_i\) và có giá trị \(V_i\).
Nếu bài đó đọc xong tại thời điểm \(C_i\), chi phí đóng góp là
\[
  C_iV_i.
\]
Cần tìm thứ tự đọc để tổng chi phí nhỏ nhất. Ngoài ra, các bài của cùng một
tác giả phải được đọc theo đúng thứ tự thời gian mà tác giả viết ra.

Ràng buộc theo tác giả làm bài toán khó hơn một bài sắp xếp thông thường. Mạch
slide không xử lý ngay trường hợp đầy đủ, mà trước hết bỏ ràng buộc để tìm
quy luật nền.

\section{Trường hợp không có ràng buộc tác giả}

Xét hai bài kề nhau trong một thứ tự, với thời gian và giá trị lần lượt là
\((T_1,V_1)\) và \((T_2,V_2)\). Nếu đổi chỗ chúng, bài thứ nhất hoàn thành
muộn thêm \(T_2\), còn bài thứ hai hoàn thành sớm hơn \(T_1\). Biến thiên của
tổng chi phí là
\[
  T_2V_1 - T_1V_2.
\]
Trong thứ tự tối ưu, đổi chỗ hai phần tử kề nhau không được làm chi phí giảm.
Do đó cần có
\[
  T_2V_1 > T_1V_2,
\]
tương đương
\[
  \frac{V_1}{T_1} > \frac{V_2}{T_2}.
\]

Vậy nếu không có ràng buộc tác giả, thứ tự tối ưu là sắp giảm dần theo
\(V_i/T_i\), với độ phức tạp \(\Oh(N\log N)\). Đây là phần ``đi sâu'' đầu
tiên: từ bài toán tổng quát rút về trường hợp tự do, rồi từ đổi chỗ cục bộ
suy ra quy luật toàn cục.

\section{Đưa ràng buộc tác giả trở lại}

Khi có nhiều bài cùng tác giả, không thể sắp từng bài riêng lẻ theo tỉ số
\(V_i/T_i\), vì thứ tự nội bộ của tác giả đã cố định. Ta lại đi từ trường hợp
nhỏ hơn: hai tác giả, trong đó một phía chỉ có một bài độc lập. Câu hỏi là
bài độc lập nên được chèn vào đâu trong dãy bài của tác giả kia.

Điều kiện đổi chỗ vẫn nói rằng tỉ số của bài được chèn phải nằm giữa tỉ số
của các phần tử lân cận. Nhưng nếu chỉ nhìn đại số, có thể có nhiều vị trí và
khó thấy cấu trúc tổng quát. Đây là lúc cần ``độ rộng'': chuyển cặp
\((V_i,T_i)\) thành vectơ trên mặt phẳng.

Tỉ số \(V_i/T_i\) chính là hệ số góc của vectơ \((T_i,V_i)\) hoặc tương đương
với hướng của cặp đại lượng. Khi ghép nhiều bài liên tiếp của cùng tác giả,
ta cộng các vectơ. Vấn đề chèn bài độc lập trở thành xét các hướng của các
đoạn trên đường gấp khúc tạo bởi dãy vectơ.

\section{Điểm lồi, điểm lõm và bao lồi}

Slide dùng hình học để phân biệt vị trí có thể giữ trong nghiệm tối ưu và vị
trí nên bị gộp. Nếu đường gấp khúc có một điểm lõm theo hướng xét, việc chèn
một bài hoặc một đoạn khác quanh điểm đó không thể tạo ra thứ tự tốt nhất.
Hai đoạn kề điểm lõm nên được gộp lại thành một đoạn lớn hơn.

Lặp thao tác gộp các điểm lõm sẽ để lại đường bao lồi. Mỗi cạnh của bao lồi
tương ứng với một đoạn liên tiếp các bài của cùng tác giả. Các đoạn này là
đơn vị có thể đem ra sắp xếp giữa các tác giả, còn thứ tự bên trong mỗi đoạn
vẫn giữ nguyên.

Cách nhìn này nối hai phần của bài nói:
\begin{itemize}
  \item chiều sâu: xuất phát từ đổi chỗ hai phần tử kề nhau và giữ mạch suy
  luận khi quay lại bài toán có ràng buộc;
  \item độ rộng: đổi biểu diễn từ tỉ số đại số sang hình học, rồi nhận ra vai
  trò của bao lồi.
\end{itemize}

\section{Thuật toán từ slide}

Thuật toán được trình bày ở mức ý tưởng như sau:
\begin{enumerate}
  \item với mỗi tác giả, giữ nguyên thứ tự các bài của tác giả đó;
  \item xem mỗi bài là một vectơ mang hai đại lượng thời gian và giá trị;
  \item trên dãy vectơ của từng tác giả, gộp các đoạn tạo điểm lõm để thu được
  các cạnh của bao lồi;
  \item mỗi cạnh bao lồi trở thành một khối bài liên tiếp;
  \item sắp tất cả các khối theo hệ số góc giảm dần, tương tự trường hợp
  không ràng buộc;
  \item bung các khối theo thứ tự nội bộ ban đầu để có thứ tự đọc cuối cùng.
\end{enumerate}

Nếu các thao tác bao lồi trên từng tác giả được thực hiện bằng ngăn xếp, tổng
số lần thêm và gộp là tuyến tính theo số bài. Bước sắp các khối chi phối thời
gian, nên độ phức tạp tổng quát vẫn ở mức \(\Oh(N\log N)\) như mốc xuất hiện
trong slide.

\section{Thông điệp phương pháp}

Ví dụ không nhằm giới thiệu bao lồi như một mẹo riêng lẻ. Nó minh họa cách
phân tích có trật tự:
\begin{itemize}
  \item giảm bài toán xuống trường hợp tự do để tìm quy luật nền;
  \item dùng đổi chỗ cục bộ để chứng minh quy luật chứ không đoán bằng trực
  giác;
  \item quay lại ràng buộc và xét trường hợp nhỏ nhất còn mang bản chất mới;
  \item đổi góc nhìn khi công thức đại số chưa đủ bộc lộ cấu trúc;
  \item tổng hợp kết luận thành thuật toán.
\end{itemize}

``Sâu'' giúp ta không bỏ qua mối quan hệ cốt lõi giữa \(T_i\) và \(V_i\).
``Rộng'' giúp ta không bị kẹt trong một biểu thức tỉ số, mà nhìn thấy hình học
đằng sau nó. Hai hướng này bổ sung cho nhau: đi sâu mà thiếu rộng dễ bế tắc;
đi rộng mà thiếu sâu dễ thành thử sai rời rạc.

\section{Kết luận}

Bản trình chiếu khép lại bằng lời nhắc về thói quen phân tích. Trong luyện
tập, giải được bài chưa phải điểm dừng; cần ghi lại con đường đi từ trường hợp
đặc biệt tới tổng quát, từ công thức tới hình học, từ quan sát tới chứng minh.
Khi ta biết vì sao thuật toán đúng, không chỉ biết nó chạy ra đáp án, năng lực
giải bài mới tăng lên một cách bền vững.

\end{document}
